\documentclass{article}
\usepackage{algorithmic}
\usepackage[T1]{fontenc}
\usepackage[utf8]{inputenc}
\usepackage{textcomp}
\usepackage{graphicx}
%\usepackage{esint}
\usepackage[shortlabels]{enumitem}
\usepackage[spanish]{babel}
\usepackage{subfig}

\usepackage{bbm}
\usepackage{amsthm}
\newtheorem{theorem}{Teorema}
\usepackage{amsthm}
\usepackage{gensymb}
\usepackage{appendix}
\newcommand\myeq{\stackrel{\mathclap{\normalfont\mbox{*}}}{=}}
\usepackage[top=1.3cm, bottom=2.0cm, outer=2.5cm, inner=2.5cm, heightrounded,
marginparwidth=2cm, marginparsep=0.4cm, margin=2.5cm]{geometry}
\usepackage{mathtools, amssymb, bm, commath, amsmath, esvect, siunitx, physics, blindtext, hyperref}


\usepackage{changes}
\numberwithin{equation}{section}

 \usepackage{todonotes}
\usepackage{marginnote}
\let\marginpar\marginnote
\newcommand{\cometex}[1]{\textrm{\todo{#1}}}
\begin{document}
	\title{Estudio de modelos matemáticos estructurados para la transmisión del dengue}
	\date{}
	\author{Teo Demergasso}

\maketitle

\section{Modelo Matemático}
El modelo presentado es una generalización del modelo planteado por de Araújo et al.\ (2023) en \emph{Applying a multi-strain dengue model to epidemics data} para un número $n$ de parches (o nodos) espaciales, con el objetivo de examinar el efecto del movimiento del hospedador en la dinámica de transmisión del dengue. Consideramos $c$ serotipos circulantes; siguiendo el trabajo original fijamos $c=4$ en la exposición, aunque la implementación numérica que acompaña a este documento (sección~\ref{sec:codigo}) admite $c$ arbitrario y usa $c=2$ por defecto. Modelamos las infecciones primarias y secundarias por dengue. Para este fin, la población humana del parche $l$, $l=1,\dots,n$, se divide en: susceptibles a los cuatro serotipos, $\bar S_l$; infectados primariamente por el serotipo $i$, $\bar I_{li}$, $i=1,\dots,4$; recuperados de la infección primaria por el serotipo $i$ pero susceptibles a los serotipos $j\neq i$, $\bar Z_{li}$; infectados primariamente por el serotipo $i$ y secundariamente por el serotipo $j\neq i$, $\bar Y_{lij}$\footnote{Convención de índices de $Y$: el \textbf{primer} índice de serotipo es la infección \textbf{primaria} y el \textbf{segundo} la \textbf{secundaria}. Es la convención de la implementación numérica. El trabajo de de Araújo et al.\ usa la opuesta ($Y_{ki}$ = primaria por $k$, secundaria por $i$, sumando sobre el primer índice); un lector que compare ambos textos debe tenerlo presente.}; y recuperados de ambas infecciones, $\bar Z_l$, que no depende del serotipo porque un individuo con dos infecciones es inmune a todos los serotipos considerados. La población de mosquitos del parche $l$ se divide en susceptibles, $\bar V_{sl}$, e infectados por el serotipo $i$, $\bar V_{li}$, $i=1,\dots,4$ (suponemos que, debido a su corta vida, los mosquitos solo pueden infectarse con dengue una vez). El total de humanos del parche $l$ capaces de transmitir el serotipo $i$ es $\bar I_{li}+\sum_{k\neq i}\bar Y_{lki}$: los infectados primarios por $i$ más los infectados secundarios por $i$, cualquiera sea su serotipo primario $k$.

Un individuo susceptible del parche $l$ se infecta con el serotipo $i$ al ser picado por un mosquito infectado del parche $m$ a una tasa $\lambda_{lim}\bar V_{mi}$, con $i=1,\dots,4$ y $l,m=1,\dots,n$. Los individuos que han sido infectados con el serotipo $i$ pueden tener una segunda infección con un serotipo diferente $j$ a una tasa $\sigma_{ij}\lambda_{ljm}\bar V_{mj}$, con $i,j=1,\dots,4$, $i\neq j$. El factor $\sigma_{ij}$ es el efecto cruzado de la infección primaria por $i$ sobre la susceptibilidad al serotipo $j$; depende solo del par de serotipos y no del parche. Un individuo infectado por un serotipo se vuelve permanentemente inmune a este serotipo, pero susceptible a los demás, lo que lleva a casos secundarios. Si $\sigma_{ij}=1$, la infección por el serotipo $i$ no altera la susceptibilidad a la infección por el serotipo $j$. Si $0\leq\sigma_{ij}<1$, el serotipo $i$ confiere inmunidad parcial o total al serotipo $j$. Si $\sigma_{ij}>1$, la infección por el serotipo $i$ aumenta la susceptibilidad al serotipo $j$ (\emph{antibody-dependent enhancement}, ADE). Por lo tanto, en nuestro modelo $\sigma_{ij}$ puede asumir valores menores o mayores que uno, ya que representa el resultado neto del efecto cruzado sobre la intensidad de la tasa de infección. Adoptamos además la convención
\begin{equation}\label{eq:sigma-diag}
\sigma_{ii}=0,\qquad i=1,\dots,4,
\end{equation}
que expresa que no hay reinfección por el mismo serotipo; con ella, las sumas sobre el segundo índice de $\sigma$ pueden escribirse sobre todos los serotipos sin la restricción $j\neq i$, y $Y_{lii}\equiv 0$. (Es exactamente lo que hace la implementación al multiplicar $\sigma$ por una máscara con ceros en la diagonal.)

Debido a la corta vida de los mosquitos, no pueden reinfectarse con otro serotipo. Además, aunque los humanos pueden contraer dengue hemorrágico (la forma más grave de la enfermedad), la mortalidad por esta enfermedad es baja, por lo que podemos ignorarla. Asimismo, asumimos las mismas tasas de mortalidad y recuperación humanas para todos los serotipos, ya que no encontramos evidencia que sugiera alguna diferencia. Según estos supuestos, podemos suponer que las poblaciones humana y de mosquitos de cada parche, denotadas por $N_{hl}$ y $N_{vl}$ respectivamente, son constantes, con tasas de mortalidad natural dadas por $\mu$ y $\nu$.

\subsection{Adimensionalización}\label{sec:adimensional}

Para simplificar el modelo tomamos proporciones en ambas poblaciones, \emph{normalizando por parche}:
\begin{equation}\label{eq:proporciones}
S_l=\frac{\bar S_l}{N_{hl}},\quad
I_{li}=\frac{\bar I_{li}}{N_{hl}},\quad
Z_{li}=\frac{\bar Z_{li}}{N_{hl}},\quad
Y_{lij}=\frac{\bar Y_{lij}}{N_{hl}},\quad
Z_{l}=\frac{\bar Z_{l}}{N_{hl}},\quad
V_{li}=\frac{\bar V_{li}}{N_{vl}},\quad
V_{sl}=\frac{\bar V_{sl}}{N_{vl}},
\end{equation}
para $l=1,\dots,n$ e $i,j=1,\dots,4$. Como cada parche tiene población constante, las proporciones satisfacen
\begin{equation}\label{eq:cierre}
Z_l = 1 - S_l - \sum_{i=1}^4\Big(I_{li}+Z_{li}+\sum_{j\neq i} Y_{lij}\Big),
\qquad
V_{sl} = 1-\sum_{i=1}^4 V_{li},
\end{equation}
de modo que $Z_l$ y $V_{sl}$ no necesitan ecuación propia: quedan determinadas por las demás variables del mismo parche.

Al normalizar por parche aparece una consecuencia que conviene explicitar. Escribamos el contagio humano en números absolutos con una tasa de contacto por mosquito: un mosquito infectado del parche $m$ infecta humanos del parche $l$ a una tasa $\tilde\lambda_{lim}$ por mosquito y por fracción de humanos susceptibles, es decir,
\[
\frac{d\bar S_l}{dt}=\mu N_{hl}-\bar S_l\sum_{m=1}^n\sum_{i=1}^4\tilde\lambda_{lim}\,\frac{\bar V_{mi}}{N_{hl}}-\mu\bar S_l .
\]
Dividiendo por $N_{hl}$ y usando \eqref{eq:proporciones},
\[
\frac{dS_l}{dt}=\mu-S_l\sum_{m=1}^n\sum_{i=1}^4\tilde\lambda_{lim}\,\frac{N_{vm}}{N_{hl}}\,V_{mi}-\mu S_l ,
\]
así que en los términos de contagio cruzado entre parches aparecen las \emph{razones de tamaños poblacionales} $N_{vm}/N_{hl}$: un parche con muchos mosquitos por humano recibe una fuerza de infección proporcionalmente mayor. Definimos entonces
\begin{equation}\label{eq:lambda-razon}
\lambda_{lim}=\tilde\lambda_{lim}\,\frac{N_{vm}}{N_{hl}},
\end{equation}
que es el coeficiente que aparece en el sistema en proporciones. Del lado del mosquito, la infección de un mosquito del parche $l$ depende de la \emph{fracción} de humanos infecciosos del parche $m$, $(\bar I_{mi}+\sum_{k\neq i}\bar Y_{mki})/N_{hm}=I_{mi}+\sum_{k\neq i}Y_{mki}$, de modo que al dividir por $N_{vl}$ no aparece ninguna razón poblacional: $\delta_{lim}$ es directamente la tasa por fracción de humanos infecciosos. Por la misma razón, con $S_l$ proporción el término de natalidad en $\dot S_l$ es $\mu$ y no $\mu N_{hl}$.

En la implementación numérica que acompaña a este documento todas las poblaciones se toman iguales, $N_{hl}=N_{vl}=1$ para todo $l$ (el código las lleva explícitamente en un vector \texttt{N\_pop} de unos), por lo que las razones \eqref{eq:lambda-razon} valen $1$, $\lambda_{lim}=\tilde\lambda_{lim}$, y el término de natalidad se escribe $\mu N_l$ con $N_l=1$. Por eso esas razones no aparecen en el código; si se quisieran parches con poblaciones distintas habría que introducirlas en $\lambda$.

De acuerdo con las suposiciones anteriores, obtenemos el siguiente sistema de ecuaciones diferenciales ordinarias con 25 ecuaciones por cada nodo ($1+4+4+12+4$; las $Y_{lii}$ son idénticamente nulas), donde $\gamma$ es la tasa de recuperación humana, común a todos los serotipos y parches, $\mu$ la tasa de natalidad y mortalidad humana y $\nu$ la mortalidad del mosquito:

\begin{equation}\label{eq:sistema}
\begin{aligned}
    \frac{d S_l}{d t} & =\mu-S_l\sum_{m=1}^n\sum_{i=1}^4 \lambda_{lim} V_{mi} -\mu S_l, \\
    \frac{d I_{li}}{d t} & =S_l\sum_{m=1}^n\lambda_{lim} V_{mi} -\left(\gamma+\mu\right) I_{li}, \\
    \frac{d Z_{li}}{d t} & =\gamma I_{li}-Z_{li}\sum_{j=1}^4\sigma_{ij}\sum_{m=1}^n  \lambda_{ljm}V_{mj}-\mu Z_{li}, \\
    \frac{d Y_{lij}}{d t} & =\sigma_{ij}Z_{li}\sum_{m=1}^n\lambda_{ljm}V_{mj}-\left(\mu+\gamma\right) Y_{lij}, \\
    \frac{d V_{li}}{d t} & =\left(1-\sum_{k=1}^4 V_{lk}\right)\sum_{m=1}^n\delta_{lim}\left(I_{mi}+\sum_{k\neq i} Y_{mki}\right)-\nu V_{li},
\end{aligned}
\end{equation}
para $l=1,\dots,n$ e $i,j=1,\dots,4$. En la ecuación de $Z_{li}$ la suma sobre $j$ recorre en efecto solo los $j\neq i$ por \eqref{eq:sigma-diag}; en la de $V_{li}$, $\sum_{k\neq i}Y_{mki}$ son los humanos del parche $m$ con infección secundaria por $i$, cualquiera sea su serotipo primario $k$. Las matrices de acoplamiento son $\lambda_{lim}$ (fuerza de infección sobre los humanos del parche $l$ ejercida por los mosquitos del parche $m$, serotipo $i$) y $\delta_{lim}$ (infección de los mosquitos del parche $l$ por los humanos del parche $m$, serotipo $i$); su forma concreta para una grilla se da en la sección~\ref{sec:experimentos}.

\subsection{Identidad de conservación}\label{sec:conservacion}

El sistema \eqref{eq:sistema} no lleva una ecuación para los recuperados de ambas infecciones, $Z_l$: esa masa sale del sistema rastreado por el término $-\gamma Y_{lij}$ y no vuelve (no hay reinfección triple). Conviene hacer explícito qué se conserva. Sea
\begin{equation}\label{eq:T}
T_l = S_l+\sum_{i=1}^4 I_{li}+\sum_{i=1}^4 Z_{li}+\sum_{i,j=1}^4 Y_{lij}
\end{equation}
la proporción de humanos del parche $l$ que el sistema rastrea. Sumando las cuatro ecuaciones humanas de \eqref{eq:sistema}, los términos de transferencia entre compartimentos se cancelan por pares: la incidencia primaria $S_l\sum_m\lambda_{lim}V_{mi}$ sale de $S_l$ y entra en $I_{li}$; la recuperación $\gamma I_{li}$ sale de $I_{li}$ y entra en $Z_{li}$; la incidencia secundaria $\sigma_{ij}Z_{li}\sum_m\lambda_{ljm}V_{mj}$ sale de $Z_{li}$ y entra en $Y_{lij}$. Solo sobreviven la natalidad, la mortalidad y la salida por recuperación de la infección secundaria:
\begin{equation}\label{eq:conservacion}
\frac{d T_l}{dt}=\mu N_l-\mu\,T_l-\gamma\sum_{i,j=1}^4 Y_{lij},
\end{equation}
donde escribimos $N_l$ por la población normalizada del parche ($N_l=1$ en proporciones; la implementación la lleva explícita, sección~\ref{sec:adimensional}). Si se introduce el compartimento faltante con su ecuación natural,
\begin{equation}\label{eq:Zl}
\frac{dZ_l}{dt}=\gamma\sum_{i,j=1}^4 Y_{lij}-\mu Z_l,
\end{equation}
entonces $\frac{d}{dt}(T_l+Z_l)=\mu\,(N_l-T_l-Z_l)$, de modo que
\begin{equation}\label{eq:invariante}
T_l(t)+Z_l(t)=N_l\qquad\text{para todo }t\ge 0
\end{equation}
es un invariante \emph{exacto} del sistema ampliado siempre que se inicialice $Z_l(0)=N_l-T_l(0)$. Esto es lo que justifica la relación de cierre \eqref{eq:cierre}: la implementación no integra $Z_l$, lo obtiene por diferencia.

La identidad \eqref{eq:conservacion} vale también refinada por serotipo primario. Sumando la ecuación de $Z_{li}$ con las de $Y_{lij}$ sobre $j$, la incidencia secundaria se cancela y queda
\begin{equation}\label{eq:conservacion-cepa}
\frac{d}{dt}\Big(Z_{li}+\sum_{j=1}^4 Y_{lij}\Big)=\gamma I_{li}-\mu Z_{li}-(\mu+\gamma)\sum_{j=1}^4 Y_{lij},\qquad i=1,\dots,4,
\end{equation}
que dice que la masa que sale de $Z_{li}$ por infección secundaria entra en $Y_{li\,\cdot}$, con el mismo serotipo primario, y no en otro lugar.

Estas identidades son el contenido del test T2 de la suite numérica que acompaña al documento (\texttt{tests/test\_conservacion.py}): se evalúa el lado derecho de \eqref{eq:sistema} en estados admisibles aleatorios y se verifican \eqref{eq:conservacion} y \eqref{eq:conservacion-cepa} con tolerancia relativa $10^{-12}$, sin integrar. Un error de índice que mande masa a un compartimento equivocado rompe \eqref{eq:conservacion}; uno que la mande al serotipo equivocado rompe \eqref{eq:conservacion-cepa}. El test T3 verifica la versión integrada, \eqref{eq:invariante}, sobre una trayectoria completa, integrando $Z_l$ como estado extra.

\section{Equilibrio Endémico}
Queremos saber bajo que condiciones existe solución al sistema 
\begin{align*}
    0 & =\mu-S_l\sum_{m=1}^n\sum_{i=1}^4 \lambda_{lim} V_{mi} -\mu S_l, \\
    0 & =S_l\sum_{m=1}^n\lambda_{lim} V_{mi} -\left(\gamma+\mu\right) I_{li}, \\
    0 & =\gamma I_{li}-Z_{li}\sum_{j=1}^4\sigma_{ij}\sum_{m=1}^n  \lambda_{ljm}V_{mj}-\mu Z_{li}, \\
    0 & =\sigma_{ij}Z_{li}\sum_{m=1}^n\lambda_{ljm}V_{mj}-\left(\mu+\gamma\right) Y_{lij}, \\
    0 & =\left(1-\sum_{k=1}^4 V_{lk}\right)\sum_{m=1}^n\delta_{lim}\left(I_{mi}+\sum_{k\neq i} Y_{mki}\right)-\nu V_{li},
\end{align*}
con $l=1,\dots,n$ e $i,j=1,\dots,4$, es decir, el sistema \eqref{eq:sistema} con los miembros izquierdos anulados.
Para simplificar el problema, consideremos primero el caso en el que hay un solo nodo ($n=1$; se omite el índice de parche):
\begin{align*}
    0 & =\mu-S\sum_{i=1}^4 \lambda_{i} V_{i} -\mu S, \\
    0 & =S\lambda_{i} V_{i} -\left(\gamma+\mu\right) I_{i}, \\
    0 & =\gamma I_{i}-Z_{i}\sum_{j=1}^4\sigma_{ij}  \lambda_{j}V_{j}-\mu Z_{i}, \\
    0 & =\sigma_{ij}Z_{i}\lambda_{j}V_{j}-\left(\mu+\gamma\right) Y_{ij}, \\
    0 & =\left(1-\sum_{k=1}^4 V_{k}\right)\delta_{i}\left(I_{i}+\sum_{k\neq i} Y_{ki}\right)-\nu V_{i},
\end{align*}
con $i,j=1,\dots,4$.

\paragraph{Versión preliminar}
Específicamente, debemos buscar condiciones que aseguran la ocurrencia de un equilibrio endémico, es decir, un equilibrio diferente del equilibrio libre de enfermedad. 

El sistema de ecuaciones se puede escribir de la forma $\Phi(x) = 0$ donde $x$ representa el vector de variables de estado de las poblaciones humanas y de mosquitos. Sea $\Omega$ el dominio biológico, este es un subconjunto abierto y acotado en $\mathbb{R}^{25}$. La clausura $\text{Cl}\text{ }\Omega$ representa la región biológicamente factible, la cual se define por las siguientes restricciones:
\begin{itemize}
    \item Condición de no negatividad: Todas las variables de los compartimentos humanos (susceptibles $S$, infectados primarios $I_i$, recuperados primarios $Z_i$, infectados secundarios $Y_{k i}$ ) y de los vectores (infectados $V_i$) deben ser mayores o iguales a cero $S \geq 0, \quad I_i \geq 0, \quad Z_i \geq 0, \quad Y_{k i} \geq 0, \quad V_i \geq 0$ para $i, k \in\{1,2,3,4\}$, $k \neq i$.

    \item Condición de conservación de la población humana: La proporción de humanos totalmente recuperados e inmunes a todos los serotipos se define como $Z$. Puesto que $Z$ no puede ser negativa ($Z \geq 0$), la suma del resto de los compartimentos humanos está acotada por $S+ \sum_{i=1}^4\left(I_i+Z_i+\sum_{k=1, k \neq i}^4 Y_{i k}\right) \leq 1$.
    
    \item Condición de conservación de la población de mosquitos: De manera análoga, la proporción de mosquitos susceptibles se define como $V_s=1-\sum_{i=1}^4 V_i$. Como la población susceptible no puede ser negativa $\left(V_s \geq 0\right)$, los mosquitos infectados están acotados por: $\sum_{i=1}^4 V_i \leq 1$.
\end{itemize}

El Teorema de Existencia de Puntos Singulares (teorema 4.2 de ''Geometrical Methods of Nonlinear Analysis - M. A. Krasnosel'Skii'') establece que si un campo vectorial continuo $\Phi$ no tiene singularidades en la frontera $\partial \Omega$ y su "rotación" $\gamma(\Phi, \Omega)$ es distinta de cero, entonces $\Phi$ tiene al menos un punto singular en el interior de $\Omega$. El Teorema 4.3 a su vez establece que, para un conjunto abierto y acotado, $\Omega$ donde el campo no tiene singularidades en su frontera $\partial \Omega$ la rotación total $\gamma(\Phi,\Omega)$ es igual a la suma de los índices de sus puntos singulares aislados: $\gamma(\Phi,\Omega) = \sum_{i=1}^s \text{ind}(x_i, \Phi)$. Sin embargo, el equilibrio libre de enfermedad se encuentra explícitamente en la frontera de este dominio (es decir, en $\partial \Omega$), puesto que allí todas las variables infecciosas $\left(I_i, Y_{k i}, V_i\right)$ son nulas; por lo tanto, debemos considerar un dominio $\Omega'$ que sea una extensión del dominio biológico hacia el espacio de coordenadas negativas. Al expandir ligeramente el dominio, el equilibrio libre de enfermedad pasa a ser un punto estrictamente interior, y la frontera $\partial \Omega'$ no posee puntos singulares. 

Sea $x_0$ el equilibrio libre de enfermedad y $\Phi'(x)$ la matriz jacobiana de $\Phi$, el teorema 6.3 establece que si $\text{det}\text{ }\Phi'(x_0) \not= 0$ entonces $x_0$ es un \textit{punto singular aislado} (epidemiológicamente esto ocurre cuando $\mathcal{R}_0\not= 1$) y $\text{ind}(x_0,\Phi) = (-1)^{-\beta}$ con $\beta$ la suma de las multiplicidades de los autovalores reales negativos de la matriz Jacobiana en $x_0$. Cuando $\mathcal{R}_0$ pasa de ser menor a 1 a ser mayor a 1, un valor propio cruza el eje imaginario y se vuelve positivo, cambiando la paridad de $\beta$  Esto significa que el índice de $x_0$ cambia de signo (por ejemplo, de -1 a +1 o viceversa). Si asumimos que $\mathcal{R}_0>1$, $\text{ind}(x_0,\Phi)$ cambió de signo y, por ende, ya no coincide con la rotación global del dominio $\gamma(\Phi,\Omega') = \sum \text{ind}(x_i, \Phi)$; esto significa que debe existir forzosamente al menos un punto singular adicional $x^\star$ dentro del dominio que compense esta diferencia en la suma. Dado que los flujos de las ecuaciones diferenciales del modelo apuntan hacia adentro del dominio biológico (las poblaciones no pueden ser negativas), este nuevo punto singular no puede estar en la región físicamente imposible (negativa); debe estar en el interior del dominio biológicamente factible. Este punto $x^\star$ debe ser un equilibrio endémico. 

\paragraph{Versión más formal}
\begin{theorem}[Existencia de equilibrio endémico]\label{thm:endemic_existence}
Considérese el sistema estacionario asociado al modelo mononodal multicepa de transmisión de dengue,
\[
\Phi(x)=0, \qquad \Phi:\mathbb{R}^{25}\to\mathbb{R}^{25},
\]
donde $x$ representa el vector de variables de estado humanas y vectoriales. 
Sea $\Omega \subset \mathbb{R}^{25}$ el dominio biológicamente factible determinado por las restricciones de no negatividad y conservación de población.

Supónganse las siguientes hipótesis:

\begin{enumerate}
\item[(H1)] $\Phi$ es continuamente diferenciable en un abierto que contiene $\overline{\Omega}$.

\item[(H2)] El dominio biológico $\Omega$ es positivamente invariante para el sistema dinámico asociado, es decir, el campo vectorial apunta hacia el interior en la frontera biológica.

\item[(H3)] Existe un abierto acotado $\Omega' \subset \mathbb{R}^{25}$ tal que 
$\overline{\Omega} \subset \Omega'$ y 
\[
\Phi(x) \neq 0 \quad \text{para todo } x \in \partial \Omega'.
\]

\item[(H4)] El equilibrio libre de enfermedad $x_0 \in \Omega$ satisface
\[
\det \Phi'(x_0) \neq 0.
\]
En particular, el número básico de reproducción $\mathcal R_0$ está bien definido y $\mathcal R_0 \neq 1$.

\item[(H5)] Cuando $\mathcal R_0$ cruza el valor crítico $1$, exactamente un autovalor real de $\Phi'(x_0)$ atraviesa transversalmente el eje imaginario.
\end{enumerate}

Entonces:

\begin{enumerate}
\item[(i)] El grado (rotación) $\gamma(\Phi,\Omega')$ está bien definido y es constante bajo variaciones continuas de los parámetros que no introduzcan ceros en $\partial \Omega'$.

\item[(ii)] Si $\mathcal R_0 < 1$, el equilibrio libre de enfermedad $x_0$ es el único punto singular de $\Phi$ en $\Omega'$ y
\[
\mathrm{ind}(x_0,\Phi)=+1.
\]

\item[(iii)] Si $\mathcal R_0 > 1$, existe al menos un punto singular adicional
\[
x^\ast \in \Omega, \qquad x^\ast \neq x_0.
\]
En particular, el sistema admite al menos un equilibrio endémico.
\end{enumerate}
\end{theorem}

\begin{proof}
La demostración se organiza en cuatro pasos, siguiendo la teoría de rotación de campos vectoriales en dimensión finita.

\medskip
\noindent
\textbf{Paso 1: Definición del grado en $\Omega'$.}

Por (H1), $\Phi$ es continuamente diferenciable en un abierto que contiene
$\overline{\Omega'}$. Por (H3), 
\[
\Phi(x) \neq 0 \quad \text{para todo } x \in \partial \Omega'.
\]
Por lo tanto, la rotación (grado topológico) $\gamma(\Phi,\Omega')$ está bien definida.

Además, por el Teorema 4.2 de Krasnosel’skii, si 
$\gamma(\Phi,\Omega') \neq 0$, entonces $\Phi$ posee al menos un punto singular
en el interior de $\Omega'$.

\medskip
\noindent
\textbf{Paso 2: Aditividad del grado e índice de puntos singulares aislados.}

Por (H4), el equilibrio libre de enfermedad $x_0$ satisface
$\det \Phi'(x_0)\neq 0$. En consecuencia, por el Teorema 6.3,
$x_0$ es un punto singular aislado y su índice viene dado por
\[
\mathrm{ind}(x_0,\Phi)
= \mathrm{sign}\big(\det \Phi'(x_0)\big).
\]

Si los puntos singulares en $\Omega'$ son aislados, el Teorema 4.3 establece que
\[
\gamma(\Phi,\Omega')
=
\sum_{x_i \in \Omega'} \mathrm{ind}(x_i,\Phi).
\]

\medskip
\noindent
\textbf{Paso 3: Caso $\mathcal R_0 < 1$.}

Cuando $\mathcal R_0 < 1$, la linealización infecciosa en $x_0$
es estable. En particular, todos los autovalores de $\Phi'(x_0)$
tienen parte real negativa y, por consiguiente,
\[
\det \Phi'(x_0) > 0.
\]
Luego
\[
\mathrm{ind}(x_0,\Phi)=+1.
\]

Por continuidad de $\Phi$ y ausencia de ceros en la frontera
(H3), el grado coincide con el índice del único punto singular
en $\Omega'$, es decir,
\[
\gamma(\Phi,\Omega')=1.
\]

\medskip
\noindent
\textbf{Paso 4: Caso $\mathcal R_0 > 1$.}

Por (H5), cuando $\mathcal R_0$ cruza el valor crítico $1$,
exactamente un autovalor real de $\Phi'(x_0)$ atraviesa
transversalmente el eje imaginario.
Por lo tanto, al pasar de $\mathcal R_0<1$ a $\mathcal R_0>1$,
la paridad del número de autovalores con parte real positiva cambia,
y en consecuencia cambia el signo de $\det \Phi'(x_0)$.

De este modo,
\[
\mathrm{ind}(x_0,\Phi)=-1.
\]

Sin embargo, el grado $\gamma(\Phi,\Omega')$ es invariante
bajo variaciones continuas de parámetros que no introduzcan ceros
en $\partial \Omega'$.
Como la condición (H3) se mantiene, el grado sigue siendo
\[
\gamma(\Phi,\Omega')=1.
\]

Aplicando la fórmula aditiva del grado,
\[
\gamma(\Phi,\Omega')
=
\mathrm{ind}(x_0,\Phi)
+
\sum_{x_i \neq x_0} \mathrm{ind}(x_i,\Phi),
\]
se obtiene
\[
1
=
(-1)
+
\sum_{x_i \neq x_0} \mathrm{ind}(x_i,\Phi),
\]
lo que implica
\[
\sum_{x_i \neq x_0} \mathrm{ind}(x_i,\Phi)=2.
\]

En particular, existe al menos un punto singular adicional
$x^\ast \in \Omega'$ distinto de $x_0$.

Finalmente, por (H2), el dominio biológico $\Omega$
es positivamente invariante. El campo vectorial apunta hacia el interior
en la frontera biológica, por lo que ningún punto singular puede ubicarse
en la región no físicamente admisible.
Por consiguiente,
\[
x^\ast \in \Omega.
\]

Este punto singular corresponde a un equilibrio endémico.

\end{proof}

\end{document}